% 原书物理页=印刷页26--33；4.1 全节（4.1.1--4.1.23，其中 4.1.20--4.1.23 在印刷页 33 上部、4.2 节标题之前）.
\chapter{线性表示}\label{chap:fg-4}
\section{定义与例子}\label{sec:fg-4-1}

设 $G$ 是一个群，$K$ 是一个交换域.

\begin{definition}\label{def:fg-4-1-1}
以 $K$ 为系数的 $G$ 的一个\newterm{fg-representation}{（线性）表示}{(linear) representation}是一个对 $(\rho,V)$，由一个 $K$-向量空间 $V$ 与一个群同态 $\rho:G\longrightarrow\mathrm{GL}_K(V)$ 组成.
\end{definition}

我们也说 $\rho$ 是一个以 $V$ 为底空间的表示. 注意到，给定 $\rho$ 等价于给定 $G$ 在 $V$ 上的一个与 $V$ 的 $K$-向量空间结构相容的左作用.

人们常犯一个情有可原的记号滥用，即把表示 $(\rho,V)$ 简单地记作 $\rho$，有时甚至更不那么说得过去地把它记作 $V$\ldots{} 此时我们把本应记作 $\rho(g)(v)$ 的东西（即 $V$ 的向量 $v$ 在 $G$ 的元素 $g$ 作用下、经由 $\rho$ 的像）简单地记作 $g.v$. 于是对一切 $g,g'\in G$、$v,v'\in V$、$\lambda\in K$，有
\[
g(\lambda v+v')=\lambda gv+gv',\quad (gg')v=g(g'v),\quad 1_Gv=v.
\]
\begin{numpara}[例]\label{np:fg-4-1-2}
\begin{enumerate}[label=\roman*)]
\item \emph{平凡表示}. 对任意 $K$-向量空间 $V$，$V$ 上的\newterm{fg-trivial-representation}{平凡表示}{trivial representation} $\mathbf{1}_V$ 定义为：对一切 $g\in G$，$\mathbf{1}_V(g)=\mathrm{id}_V$. 当 $V=K$ 时，平凡表示 $\mathbf{1}_K$ 称为 $G$ 的\newterm{fg-unit-representation}{单位表示}{unit representation}（相对于域 $K$ 而言）.
\item \emph{标准表示}. 当 $G$ 本来就作为 $\mathrm{GL}(V)$ 的子群给出时（例如 $\mathrm{GL}(2,\mathbb{C})$ 中的二面体群 $D_n$），$V$ 自然就是 $G$ 的一个表示，常称为自然表示，或\newterm{fg-standard-representation}{标准表示}{standard representation}.
\item \emph{置换表示}. 设 $X$ 是一个带有 $G$ 的左作用的集合. 相伴的\newterm{fg-permutation-representation}{置换表示}{permutation representation}是以 $\{e_x,\ x\in X\}$ 为基的 $K$-向量空间 $V=KX$，其上的作用由条件 $g.e_x=e_{g.x}$ 经 $K$-线性延拓定义. 于是 $KX$ 中的元素 $e$ 可唯一地写成 $\sum_x\lambda_x e_x$，其中 $(\lambda_x)_{x\in X}$ 是一族``几乎为零''的 $K$ 元素（即除有限多个之外全为零），并且有
\[
g.e=g.\Bigl(\sum_x\lambda_x e_x\Bigr)=\sum_x\lambda_x e_{g.x}.
\]
例：取 $G=S_n$、$X=\{1,\cdots,n\}$，便得到 $S_n$ 在 $K^n$ 上的一个表示. 同态 $S_n\longrightarrow\mathrm{GL}_n(K)$ 把一个置换映到相伴的置换矩阵. 更一般地，在基 $(e_x)_{x\in X}$ 下可以看出，给出 $g$ 的作用的矩阵是一个置换矩阵（每行、每列恰有一个非零元素且等于 1）. 置换表示之名即由此而来.
\item \emph{函数空间}. 仍从上面的 $X$ 出发，从 $X$ 到 $K$ 的函数全体构成的空间 $\mathcal{F}(X,K):=K^X$ 带有如下定义的作用：
\[
\forall x\in X,\ (g.f)(x):=f(g^{-1}x).
\]
注意，取逆这一步对于等式 $g.(h.f)=(gh).f$（即左作用）的成立是必不可少的；否则得到的将是右作用. 注意到，上面的空间 $KX$ 可以等同于 $K^X$ 的由支撑有限的函数（即在 $X$ 的有限子集之外恒为零的函数）组成的子向量空间 $\mathcal{F}_0(X,K)$. 事实上，一个同构由
\[
\begin{array}{rcl}
KX & \longrightarrow & \mathcal{F}_0(X,K)\\
e_x & \longmapsto & \delta_x
\end{array}
\]
给出，其中 $\delta_x$ 是以 $x$ 为中心的 Dirac 函数，定义为：当 $x\neq y$ 时 $\delta_x(y)=0$，而 $\delta_x(x)=1$. 由上面的公式可知 $g.\delta_x=\delta_{g.x}$. 于是上述同构与 $G$ 在 $KX$ 及 $\mathcal{F}_0(X,K)$ 上的表示相容.
\item \emph{正则表示}. 取 $X=G$，并赋予 $G$ 的左平移作用，则置换表示 $V=KG\simeq\mathcal{F}_0(G,K)$ 称为 $G$ 的\newterm{fg-regular-representation}{（左）正则表示}{(left) regular representation}，记作 $R$. 于是它定义为
\[
R(g)\Biggl(\sum_{h\in G}\lambda_h e_h\Biggr)=\sum_{h\in G}\lambda_h e_{gh}=\sum_{h\in G}\lambda_{g^{-1}h}e_h,
\]
其中对每个 $h\in G$，$\lambda_h$ 是 $K$ 的一个元素，且在 $G$ 的某个有限子集之外 $\lambda_h=0$.
\end{enumerate}
\end{numpara}

\begin{numpara}[小词典]\label{np:fg-4-1-3}
\emph{表示的\newterm{fg-degree}{维数}{degree}.}~即底 $K$-向量空间在 $K$ 上的维数.

\emph{\newterm{fg-fixed-vector}{固定向量}{fixed vector}.}~这是 $G$ 在 $V$ 上作用的一个不动点，即满足 $gv=v$ 的元素 $v\in V$. 零向量 $0\in V$ 总被 $G$ 固定. 我们常把在 $G$ 下固定的向量全体记作 $V^G$；它是 $V$ 的一个 $K$-子向量空间.

\noindent\textbf{例.}
\begin{enumerate}[label=\roman*)]
\item 在置换表示 $KX$ 中，向量 $\sum_{x\in X}e_x$ 在 $G$ 下固定.
\item 若 $G$ 是有限 $p$-群而 $K$ 的特征为 $p$，则任何表示都有非零的固定向量. 事实上，从非零的 $v\in V$ 出发，考察由向量 $gv$（$g\in G$）生成的 $\mathbb{F}_p$-向量空间 $W$. 它是 $G$ 的一个有限维 $\mathbb{F}_p$-表示，其元素个数为某个 $p^n$. 至少存在一个固定向量，即零向量；而 $G$ 在 $W$ 上作用的类公式告诉我们，固定向量的个数被 $p$ 整除.
\end{enumerate}
\emph{\newterm{fg-faithful-representation}{忠实表示}{faithful representation}.}~指同态 $\rho$ 为单射的表示 $(\rho,V)$.
\emph{\newterm{fg-subrepresentation}{子表示}{subrepresentation}.}~指 $V$ 的一个在 $G$ 作用下\newterm{fg-stable-subspace}{稳定}{stable}（人们也说：``不变的''）的 $K$-子向量空间 $W$. 这意味着对一切 $v\in W$ 与一切 $g\in G$ 都有 $gv\in W$. 严格说来，应当把它记作 $(\rho,W)\subset(\rho,V)$.

\emph{\newterm{fg-quotient-representation}{商表示}{quotient representation}.}~若 $W$ 是 $V$ 的一个稳定子空间，则商向量空间 $V/W$ 自然地带有 $G$ 的表示结构，称为 $V$ 对 $W$ 的商表示. 对每个陪集 $v+W$，简单地令 $g(v+W):=gv+W$；由于 $W$ 假定稳定，这与 $v$ 的选取无关.
\end{numpara}

\begin{definition}\label{def:fg-4-1-4}
$G$ 的两个表示之间的一个\newterm{fg-gmorphism}{表示间的同态}{morphism of representations} $(\rho,V)\longrightarrow(\pi,W)$ 是一个 $K$-线性同态 $\varphi:V\longrightarrow W$，使得对一切 $(g,v)\in G\times V$ 都有 $\varphi(g.v)=g.\varphi(v)$.
\end{definition}

我们也说 $G$-态射，或者说 $G$-等变线性映射. 照例有各种变体：\emph{单态射} = 单的 $G$-态射，\emph{满态射} = 满的 $G$-态射，\emph{同构} = 双射的 $G$-态射（此时我们说这两个表示\emph{同构}，或者说\emph{等价}），当 $V=W$ 时称\emph{自同态}，而\emph{自同构} = 双射的自同态.

此外，$G$-态射 $\varphi:V\longrightarrow W$ 的像是 $W$ 的一个子表示，其核是 $V$ 的一个子表示. 与向量空间的情形一样，$\varphi$ 此时诱导出一个 $G$-同构
\[
V/\operatorname{Ker}(\varphi)\xrightarrow{\ \sim\ }\operatorname{Im}(\varphi).
\]
\begin{numpara}[例]\label{np:fg-4-1-5}
\begin{enumerate}[label=\roman*)]
\item 由 $V$ 的固定向量 $v$ 生成的 $V$ 的子表示 $W=K.v$ 同构于 $G$ 的单位表示 $\mathbf{1}_K$.
\item 对一个 $G$-集 $X$，前面定义的线性映射 $KX\longrightarrow K^X$ 是一个 $G$-单态射，并诱导出一个 $G$-同构 $KX\xrightarrow{\ \sim\ }\mathcal{F}_0(X,K)$.
\item 对两个 $G$-集 $X$ 与 $Y$，任何 $G$-等变映射 $X\longrightarrow Y$ 经线性延拓成为一个 $G$-态射 $KX\longrightarrow KY$，并诱导出一个 $G$-态射 $K^Y\longrightarrow K^X$.
\end{enumerate}
\end{numpara}

\begin{definition}\label{def:fg-4-1-6}
称表示 $(\rho,V)$ 是\newterm{fg-irreducible-representation}{不可约的}{irreducible}，如果它没有真子表示（即异于 $\{0\}$ 与 $V$ 的子表示）.
\end{definition}

\begin{numpara}[例]\label{np:fg-4-1-7}
\begin{enumerate}[label=\roman*)]
\item \emph{若 $G$ 是有限 Abel 群，则 $G$ 在 $\mathbb{C}$ 上的不可约表示的维数为 1.} 事实上，对每个阶为 $o(g)$ 的元素 $g$，$V$ 的自同态 $\rho(g)$ 的极小多项式整除 $X^{o(g)}-1$. 于是族 $(\rho(g))$ 是一族两两交换的可对角化自同态，而我们知道这样的族是``可同时对角化''的，即存在一组由所有 $\rho(g)$ 的特征向量构成的基. 每个特征向量生成一条在 $G$ 下稳定的向量直线.
\item 二面体群 $D_n$ 在 $\mathbb{C}^2$ 上的标准表示是不可约的.
\item $S_n$ 在 $K^n$ 上的置换表示不是不可约的，因为由 $e_1+\cdots+e_n$ 生成的直线在 $S_n$ 下稳定. 更一般地，维数有限且 $\neq 1$ 的置换表示永远不是不可约的.
\item \emph{若 $G$ 是 $p$-群而 $K$ 的特征为 $p$，则唯一的不可约表示是（在同构意义下）平凡表示 $\mathbf{1}_K$.} 这事实上由上面看到的事实——任何表示都有非零的固定向量——推出.
\item 让我们把例子 i) 与 iv) 结合起来. 设 $G=\mathbb{Z}/p\mathbb{Z}$，并设 $K$ 代数闭. 那么 $G$ 的不可约表示都是 1 维的，且形如 $\rho_\zeta:\ n\in G\mapsto\zeta^n\in\mathrm{GL}_1(K)=K^\times$，其中 $\zeta$ 是一个 $p$ 次单位根. 但随 $K$ 的特征不同，情形有着巨大的差别. 若该特征 $\neq p$，则存在 $p$ 个两两不同构的不可约表示（$p$ 个 $p$ 次单位根）. 若特征等于 $p$，则只有一个（平凡表示），因为唯一的 $p$ 次单位根是 1.
\end{enumerate}
\end{numpara}

\begin{proposition}\label{prop:fg-4-1-8}
（Schur 引理）设 $V$ 与 $W$ 是群 $G$ 的两个不可约表示，$\varphi:V\longrightarrow W$ 是一个非恒零的 $G$-态射. 那么：
\begin{enumerate}[label=\roman*)]
\item $\varphi$ 是一个同构；
\item 若再有 $V=W$ 为有限维，且 $K$ 是代数闭的\footnote{译注：原文此处作 ``si $K$ est algébriquement''，疑脱 ``clos''（闭）一词.}，则存在一个非零纯量 $\lambda$ 使得 $\varphi=\lambda.\mathrm{id}_V$.
\end{enumerate}
\end{proposition}

\begin{proof}
$\varphi$ 的核（相应地，像）是 $V$（相应地 $W$）的不变子空间. 这些表示的不可约性给出第一条断言. 在第二条断言的假设下，$\varphi$ 至少有一个特征值 $\lambda\in K$，于是只需对 $G$-态射 $\varphi-\lambda\mathrm{id}_V$ 应用第一条断言.
\end{proof}

\begin{remark*}
对一个表示 $V$，$V$ 的 $G$-自同态全体是一个 $K$-代数. 若 $V$ 是有限维的，这个 $K$-代数也是有限维的. 若 $V$ 不可约，则这个 $K$-代数是一个域（可能是非交换的，也称为\newterm{fg-division-algebra}{可除代数}{division algebra}）. 代数闭域上任何有限维可除代数都同构于该域本身.
\end{remark*}

\emph{\newterm{fg-direct-sum-representations}{直和}{direct sum}.}~若 $V$ 与 $W$ 是 $G$ 的两个表示，则它们的直和是表示 $V\oplus W$，定义为 $g.(v\oplus w)=gv\oplus gw$.

\begin{definition}\label{def:fg-4-1-9}
称表示 $(\rho,V)$ 是\newterm{fg-indecomposable-representation}{不可分解的}{indecomposable}，如果它不是两个真子表示的直和；否则称它是可分解的.
\end{definition}

注意，不可分解表示未必是不可约的. 这里有一个典型的反例.

\begin{proposition}\label{prop:fg-4-1-10}
系数取在特征为 $p$ 的域 $K$ 中时，有限 $p$-群 $G$ 的正则表示 $KG$ 是不可分解的，但不是不可约的.
\end{proposition}

\begin{proof}
前面已经看到它不是不可约的，因为它有非零的固定向量，例如向量 $v_0:=\sum_g e_g$. 事实上容易看出，$KG$ 的\emph{任何}固定向量都是 $v_0$ 的倍数. 令 $v=\sum_g\lambda_g e_g$ 是这样一个向量. 对每个 $h$，$hv=v$ 当且仅当对每个 $g$ 都有 $\lambda_g=\lambda_{hg}$. 特别地，对每个 $h$ 都有 $\lambda_h=\lambda_e$，从而 $v=\lambda_ev_0$. 由此推出，$W$ 的任何子表示都含有 $v_0$ 的一个非零倍数，从而含有 $v_0$\footnote{译注：原文此处作 ``toute sous-représentation de $W$''，$W$ 未在本证明中定义，疑为 $KG$ 之误.}. 于是两个非零子表示总有非零的交，因而更不可能构成直和.
\end{proof}

这个现象实际上正是由于 $K$ 的特征整除群的阶所致，正如下面的引理所表明的. 趁此机会我们回顾：``$K$ 的特征不整除 $G$ 的阶''与``$G$ 的阶在 $K$ 中可逆''是等价的.

\begin{proposition}\label{prop:fg-4-1-11}
（Maschke 引理）设 $G$ 是有限群，$K$ 是特征不整除 $|G|$ 的域. 设 $V$ 是一个表示，$W$ 是一个非零子表示. 那么存在 $V$ 的一个子表示 $W'$，使得 $V=W\oplus W'$.
\end{proposition}

\begin{proof}
取 $W$ 在 $V$ 中的一个（向量空间意义下的）补空间 $\tilde{W}$，以及沿 $W'$ 而到 $W$ 上的线性投影 $\tilde{p}$\footnote{译注：原文此处作 ``le long de $W'$''（沿 $W'$），但 $W'$ 迟至下文才定义为 $\operatorname{Ker}(p)$；按文意，此处应沿补空间 $\tilde{W}$ 投影.}. 它是 $\operatorname{Hom}(V,W)$ 的一个元素，但一般不是 $G$-态射. 于是令
\[
p=|G|^{-1}\Biggl(\sum_{g\in G}\rho(g)\circ\tilde{p}\circ\rho(g)^{-1}\Biggr).
\]
可以验证 $p$ 仍是 $V$ 到 $W$ 上的一个投影，而这一次它在 $G$ 下不变. 因此它的核 $W'=\operatorname{Ker}(p)$ 是 $W$ 在 $V$ 中的一个补空间，且在 $G$ 下稳定.
\end{proof}

\begin{corollary}\label{cor:fg-4-1-12}
设 $G$ 是一个阶在 $K$ 中可逆的有限群.
\begin{enumerate}[label=\roman*)]
\item 任何不可分解表示都是不可约的；
\item 任何表示 $\rho$ 都分解为不可约表示的一个直和 $\rho=\rho_1\oplus\cdots\oplus\rho_k$.
\end{enumerate}
\end{corollary}

在 ii) 中把彼此同构的不可约表示归并到一起，就得到一个分解
\[
\rho=\rho_1^{a_1}\oplus\cdots\oplus\rho_l^{a_l},
\]
其中 $a_i$ 是非零整数，$\rho_i^{a_i}$ 表示 $a_i$ 份 $\rho_i$ 的直和 $\rho_i\oplus\rho_i\oplus\cdots\oplus\rho_i$，而诸 $\rho_i$ 两两不同构. 这样的分解称为\newterm{fg-isotypic-decomposition}{同型分解}{isotypic decomposition}. 下面的引理说的是：一个同型分解的直和因子 $\rho_i$ 与整数 $a_i$ 只依赖于 $\rho$（而不依赖于得到该分解所用的归纳步骤）.

\begin{lemma}\label{lemma:fg-4-1-13}
若 $\rho=\rho_1'^{a_1'}\oplus\cdots\oplus\rho_l'^{a_l'}$ 是如上形式的另一个分解，则 $l=l'$，且存在一个置换 $\sigma\in S_l$，使得对一切 $i=1,\cdots,l$ 都有 $\rho_i\simeq\rho'_{\sigma(i)}$.
\end{lemma}

\begin{definition}\label{def:fg-4-1-14}
在如上形式的分解中出现的不可约表示称为\newterm{fg-jordan-hoelder-factor}{Jordan--Hölder 因子}{Jordan-Hölder factor}或\emph{不可约子商}. 与之相伴的整数 $a_i$ 称为\newterm{fg-multiplicity}{重数}{multiplicity}.
\end{definition}

由此可见，当系数域是特征与 $|G|$ 互素的代数闭域时，有限群 $G$ 的表示论本质上就归结为对不可约表示的分类. 若特征整除 $|G|$，理论就要复杂得多，因为还必须理解不可约表示之间的``非分裂''扩张，以及更一般的不可分解表示. 事实上（超出课程范围！），对于像 $(\mathbb{Z}/p\mathbb{Z})^2$ 这样简单的群，想得到不可分解表示的一个合理分类是毫无希望的\ldots

\begin{example}\label{ex:fg-4-1-15}
考察群 $G=\mathbb{Z}/2\mathbb{Z}$ 的正则表示. 它有一组基 $(e_0,e_1)$，$G$ 的生成元 $\gamma$ 在其上的作用为 $\gamma.e_0=e_1$ 与 $\gamma.e_1=e_0$. 若 $K$ 的特征不是 2，则可以看出向量 $e_+=e_0+e_1$ 与 $e_-=e_0-e_1$ 是 $\gamma$ 的特征向量（特征值分别为 1 与 $-1$），从而生成在 $G$ 下稳定的直线 $L_+$ 与 $L_-$. 于是立即验证 $KG=L_-\oplus L_+$. 反之，若 $K$ 的特征为 2，则 $e_+=e_-$，且 $L_+=L_-$ 是 $KG$ 中唯一在 $G$ 下稳定的直线.
\end{example}

\begin{numpara}[逆步表示（或对偶表示）]\label{np:fg-4-1-16}
设 $(V,\rho)$ 是 $G$ 的一个表示，$V^*=\operatorname{Hom}(V,K)$ 是 $K$-向量空间 $V$ 的\newterm{fg-dual}{对偶}{dual}. 对 $V$ 上的任意线性形式 $f$ 与任意 $g\in G$，令 $\rho^*(g)(f)=f\circ\rho(g^{-1})$（简记作 $g.f=f\circ g^{-1}$），从而使 $V^*$ 带有 $G$ 的表示结构；于是对一切 $(g,f,v)\in G\times V^*\times V$，有
\[
\langle g.f\mid g.v\rangle:=\langle\rho^*(g)(f)\mid\rho(g)(v)\rangle=\langle f\mid v\rangle;
\]
换言之，$\rho^*(g)\in\operatorname{End}_K(V^*)$ 是 $\rho(g^{-1})\in\operatorname{End}_K(V)$ 的转置. 我们称 $\rho^*$ 是 $\rho$ 的\newterm{fg-dual-representation}{对偶表示}{dual representation}，或\emph{逆步表示}.
\end{numpara}

\begin{numpara}[内部 Hom]\label{np:fg-4-1-17}
更一般地，设 $(\rho,V)$ 与 $(\pi,W)$ 是 $G$ 的两个表示. 从 $V$ 到 $W$ 的 $K$-线性同态全体构成的空间 $\operatorname{Hom}_K(V,W)$ 自然地带有 $G$ 的表示结构：对任意 $g\in G$ 与任意 $K$-同态 $\varphi:V\longrightarrow W$，令
\[
g.\varphi:=\pi(g)\circ\varphi\circ\rho(g^{-1}).
\]
在此作用下固定的线性映射 $\varphi$ 不外是一个 $G$-态射. 于是 $G$-态射全体 $\operatorname{Hom}_K(V,W)^G$ 是一个 $K$-子空间，习惯上记作 $\operatorname{Hom}_{KG}(V,W)$. 若 $V=W$，则 $\operatorname{End}_{KG}(V)$ 是 $\operatorname{End}_K(V)$ 的一个子 $K$-代数，有时称为 $G$ 在 $V$ 中的\newterm{fg-commutant-algebra}{换位代数}{commutant algebra}.

举两个特殊情形：表示 $V^*$ 可等同于 $\operatorname{Hom}_K(V,\mathbf{1}_K)$，而表示 $W$ 可等同于 $\operatorname{Hom}_K(\mathbf{1}_K,W)$.
\end{numpara}

\begin{numpara}[张量积]\label{np:fg-4-1-18}
我们来回顾（或解释）一个在\newterm{fg-bilinear-map}{双线性}{bilinear}代数与\emph{线性}代数之间建立联系的构造. 取定两个 $K$-向量空间 $V,W$. 对一个 $K$-向量空间 $E$，有 $K$-双线性映射 $V\times W\xrightarrow{\ \theta\ }E$ 的概念：这是指这样一类映射，使得对每个 $v\in V$，映射 $w\in W\mapsto\theta(v,w)\in E$ 是 $K$-线性的，并且同样对每个 $w\in W$，映射 $v\in V\mapsto\theta(v,w)\in E$ 是 $K$-线性的. 重要的是要记住 $\theta$ 不是线性的：例如 $\theta(2x,2y)=4\theta(x,y)$. 不过，利用\newterm{fg-tensor-product}{张量积}{tensor product}的概念，可以把双线性映射的理论化归为线性映射的理论.

为此，先注意到：若 $E\xrightarrow{\ \varphi\ }E'$ 是一个 $K$-线性映射，则 $\varphi\circ\theta:V\times W\longrightarrow E'$ 仍是 $K$-双线性的.
\end{numpara}

\begin{proposition}\label{prop:fg-4-1-19}
存在一个在同构意义下唯一的向量空间，记作 $V\otimes_K W$，称为 $V$ 与 $W$ 在 $K$ 上的张量积，它带有一个 $K$-双线性映射 $\otimes:V\times W\longrightarrow V\otimes_K W$，$(v,w)\mapsto x\otimes w$\footnote{译注：原文此处作 ``$(v,w)\mapsto x\otimes w$''，疑为 $v\otimes w$ 之误.}，并且满足如下的\newterm{fg-universal-property}{泛性质}{universal property}：对任何 $K$-双线性映射 $\theta:V\times W\longrightarrow E$，存在唯一的 $K$-线性映射 $\varphi_\theta:V\otimes_K W\longrightarrow E$，使得对一切 $(v,w)\in V\times W$ 都有 $\theta(v,w)=\varphi_\theta(v\otimes w)$.
\end{proposition}

于是，如上的一个双线性映射 $\theta:V\times W\longrightarrow E$ 由唯一的一个 $K$-线性映射 $\varphi_\theta:V\otimes_K W\longrightarrow E$ 所确定，同时它也确定后者.

$V\otimes_K W$ 的 $K$-向量空间结构可能是出错之源. 正因如此，最好把下列性质明确指出来，它们是由 $(v,w)\mapsto v\otimes w$ 的双线性性假设所强加的：
\begin{itemize}[label=--]
\item 对一切 $v,v'\in V$、$w,w'\in W$，有 $(v+v')\otimes w=v\otimes w+v'\otimes w$ 以及 $v\otimes(w+w')=v\otimes w+v\otimes w'$；
\item 对一切 $\lambda\in K$，有 $(\lambda v)\otimes w=v\otimes(\lambda w)=\lambda(v\otimes w)$.
\end{itemize}
无论如何要记住：$(v+v')\otimes(w+w')\neq v\otimes w+v'\otimes w'$.

\begin{proof}
存在性一旦证明，唯一性便像任何``泛问题的解''那样随之推出. 至于存在性，可以这样做：设 $\{e_i;\ i\in I\}$、$\{\varepsilon_j;\ j\in J\}$ 分别是 $V,W$ 的基. 考虑以符号 $e_i\otimes\varepsilon_j$（$i\in I$，$j\in J$）为基的向量空间 $\mathcal{E}$. 定义一个 $K$-双线性映射
\[
\begin{array}{rcl}
\theta_{\mathcal{E}}:\ V\times W & \longrightarrow & \mathcal{E}\\
(v,w) & \longmapsto & \sum_{i,j}(v_iw_j)e_i\otimes\varepsilon_j
\end{array}
\]
其中 $v=\sum_i v_ie_i$ 与 $w=\sum_j w_j\varepsilon_j$ 分别是 $v$ 与 $w$ 在基 $(e_i)_i$、$(\varepsilon_j)_j$ 下的分解. 对陈述中那样的任何 $K$-双线性映射 $\theta$，令 $\varphi_\theta:\mathcal{E}\longrightarrow E$ 为 $\varphi_\theta(e_i\otimes\varepsilon_j)=\theta(e_i,\varepsilon_j)$. 显然 $\theta=\varphi_\theta\circ\theta_{\mathcal{E}}$. 同样显然，$\varphi_\theta$ 是唯一满足该性质的线性映射 $\mathcal{E}\longrightarrow E$，因为一个以 $V\times W$ 为源的双线性形式由它在诸偶 $(e_i,\varepsilon_j)$（$i\in I$，$j\in J$）上的取值唯一确定.
\end{proof}

由这个证明还可以得出如何构造 $V\otimes_K W$ 的基；特别地，若 $V$ 与 $W$ 是有限维的，则 $V\otimes_K W$ 也是，并且有 $\dim(V\otimes W)=\dim(V)\dim(W)$.

现在设 $V,W$ 是群 $G$ 的表示 $\rho,\sigma$ 的空间. 那么 $V\otimes_K W$ 是 $G$ 的某个表示的空间，该表示记作 $\rho\otimes\sigma$，由如下条件定义：对一切 $(g,v,w)\in G\times V\times W$，
\[
(\rho\otimes\sigma)(g)(v\otimes w):=\rho(g)(v)\otimes\sigma(g)(w)=gv\otimes gw.
\]
我们把它称为表示 $\rho$ 与 $\sigma$ 的张量积.

\begin{example}\label{ex:fg-4-1-20}
若 $\rho$ 与 $\sigma$ 不可约，则对 $\rho\otimes\sigma$ 而言这一点往往不再成立. 例如取 $G=D_n$，$\rho$ 为 $\mathbb{C}^2$ 上的标准表示，则可以证明 $\rho\otimes\rho$ 仍包含 $\rho$.
\end{example}

\begin{numpara}[张量积与内部 Hom]\label{np:fg-4-1-21}
设 $V^*$ 是 $V$ 的对偶. 由前一命题，双线性映射 $W\times V^*\longrightarrow\operatorname{Hom}_K(V,W)$，$(w,f)\longmapsto\{\varphi:\ v\longmapsto\varphi(v):=f(v)w\}$ 定义了一个典范同态
\[
\iota:\ W\otimes_K V^*\longrightarrow\operatorname{Hom}_K(V,W).
\]
\end{numpara}

\begin{lemma}\label{lemma:fg-4-1-22}
同态 $\iota$ 是单射. 此外：
\begin{enumerate}[label=\roman*)]
\item 若 $V$ 与 $W$ 是有限维的，则它是双射的；
\item 若 $V$ 与 $W$ 都是表示，且 $W\otimes_K V^*$ 与 $\operatorname{Hom}_K(V,W)$ 分别被赋予``张量积''与\newterm{fg-internal-hom}{内部 Hom}{internal Hom}表示，则它是 $G$-等变的.
\end{enumerate}
\end{lemma}

\begin{remark*}
一般地，$\iota$ 的像由那些具有有限秩的同态 $\varphi:V\longrightarrow W$ 组成，即使得 $\varphi(V)$ 是 $W$ 的一个有限维子向量空间的那些同态.
\end{remark*}

于是，若 $(\rho,V)$ 与 $(\pi,W)$ 都是有限维表示，则 $\operatorname{Hom}_K(V,W)$ 上的表示就同构于 $\pi\otimes\rho^*$.

有一个有用的特殊情形：$V=W$ 且 $\dim_K(V)$ 有限. 此时 $\iota$ 是一个同构 $V\otimes V^*\xrightarrow{\ \sim\ }\operatorname{End}_K(V)$. $\operatorname{Id}_V$ 的原像由张量 $\sum_{i\in I}e_i\otimes e_i^*$ 给出，其中 $\{e_i\}$ 是 $V$ 的任意一组基，而 $\{e_i^*\}$ 是它在 $V^*$ 中的对偶基.

\begin{lemma}\label{lemma:fg-4-1-23}
设 $\varphi$ 是 $V$ 的一个对应于张量 $\sum_\ell v_\ell\otimes f_\ell\in V\otimes V^*$ 的自同态. 那么它的迹由公式
\[
\operatorname{Tr}(\varphi)=\sum_\ell f_\ell(v_\ell)\in K
\]
给出.
\end{lemma}
